\section{Related Work}
\label{sec:related}

Recent work on decision models has measured accuracy and calibration on datasets with one correct choice~\citep{porcedda2026jev}. \citet{deusser2026evaluating} evaluate Jev on 37 datasets and report that its probabilities are calibrated on their tasks. They also find that rotating choices leaves aggregate accuracy unchanged. On yes/no questions, Jev's probabilities rank well but are poorly placed relative to a threshold of 0.5. \citet{liu2026dice} finds that Jev gives unequal probabilities to the six faces of a fair die. Both studies test Jev only. \citet{li2026judge} find that on judging tasks with one correct answer, Jev's answers with a probability of at least 0.9 are as accurate as GPT-6. Our study directly extends their work to lists with more than one correct choice. \citet{li2026beyond} compare Jev's answers to logically related yes/no and Choice questions and find differences between probabilities for events that should agree. TypeSafe states that the order of choices can affect Jev's answer \citep{typesafe2026jev}. \citet{li2026readout} find that restricting and renormalizing the original probabilities predicts aggregate accuracy after choices are removed. For generative language models, \citet{zheng2024large,pezeshkpour2024large} show that choice order changes answers to multiple-choice questions. \citet{blankenstein2026biasbusters} find that language models selecting among functionally equivalent tools can favour particular providers. \citet{wang2025leastincorrect} find that LLMs can select an offered answer even when none is correct. However, modern AI tasks require selecting more than one correct choice. Examples include selecting the tools an agent needs to complete a task, the diagnoses supported by a medical record, the passages relevant to a query in retrieval-augmented generation, and the vulnerability types present in code.

Calibration has also been studied for multi-label classification. \citet{cheng2024calibrated} examine calibration in multi-label neural networks and propose a loss function to improve probability estimation. \citet{ullah2025calibrated} evaluate calibration in extreme multi-label classifiers, distinguishing the reliability of predicted label probabilities from ranking performance. Unlike these classifiers, the decision models we study return a normalized probability distribution over candidate choices, even when multiple choices are correct.

Human disagreement provides another setting for studying probability allocation. \citet{baan-etal-2022-stop} show that conventional calibration against majority labels can be misleading when annotators disagree. They evaluate model probabilities against human judgment distributions on ChaosNLI using measures that capture label frequencies, rankings, and entropy. We extend this investigation to decision models, examining how their probability allocations compare with empirical annotator vote distributions.
